\section{Source-model transfer and comparison scope}

The retained source restrictions fix a nonempty family \(\mathfrak P\), common class constants, and selected regression and propensity representatives before stating a uniform recovery guarantee. Under \citet[Section 2, Assumptions 1--2, author PDF pp.~3--7]{Dorn2026GlobalOverlap}, each member has uniform covariates on \([-1,1]^d\) and a selected measurable propensity satisfying
\[
e_P(X)=P(A=1\mid X)\in(0,1)
\qquad P\text{-almost surely}.
\]
Writing \(\mu_P\) for the selected treated regression, the common constants \(q>3\), \(M,\sigma_{\min},C>0\), and finite overlap exponent \(\gamma_0>1\) satisfy
\[
\mathbb E_P\!\left[|Y|^q\mid X,A=1\right]\le M^q,
\qquad
\operatorname{Var}_P(\mu_P(X))\le M,
\qquad
\operatorname{Var}_P(Y\mid X,A)\ge\sigma_{\min}^2,
\]
with conditional inequalities holding almost surely, and
\[
P\{e_P(X)\le t\}\le Ct^{\gamma_0-1}
\qquad(t\in[0,1]).
\]
Conditional outcomes in both treatment arms are Gaussian with their selected regression means and common variance \(\sigma_{\min}^2\). Both regression curves belong to \(\Sigma^{\mathrm D}(\beta,L_0)\), the source smoothness class with common \(\beta,L_0>0\).

The source smoothness convention uses derivative order
\[
m=\max\{k\in\mathbb Z:k<\beta\}=\lceil\beta\rceil-1
\]
and bounds the \((\beta-m)\)-Hölder seminorm of derivatives of order \(m\) by \(L_0\), using Euclidean distance. For the intrinsic cube interpretation used in \cref{obj:prop:dorn-a3-free-corollary}, interior derivatives through order \(m\) have continuous traces on the closed cube. The source upper comparison additionally imposes the family-level local condition in \cref{obj:def:dorn-a3}, with common positive constants and the selected pointwise propensity representatives, as in \citet[Section 2, Assumption 3, author PDF pp.~6--7]{Dorn2026GlobalOverlap}. These restrictions distinguish the retained source envelope from the fixed families to which the published upper comparison applies.

\subsection{Fixed-cube smoothness completion}

For the transfer in \cref{obj:prop:dorn-a3-free-corollary}, the treated-moment restriction supplies the function bound that complements the source seminorm. The full response class in \cref{obj:def:holder} also controls the function and all lower derivatives. The following completion statement connects these two smoothness descriptions for the paper's range \(\beta>1\).

\begin{lemmav}{lem:fixed-cube-holder-completion}[Hölder completion under affine scaling]
Fix an integer \(d\ge 1\) and a real number \(\beta>1\), and set \(m=\lceil\beta\rceil-1\). There exists a constant \(K_{d,\beta}>0\), depending only on \(d\) and \(\beta\), such that the following holds for every \(M_0,L_0\in\mathbb R\) and \(u:\mathbb R^d\to\mathbb R\) satisfying:
\begin{itemize}
\item \textbf{(Nonnegative bounds.)} \(M_0\ge0\) and \(L_0\ge0\).
\item \textbf{(Intrinsic regularity.)} \(u\in C^m([-1,1]^d)\), with interior derivatives through order \(m\) admitting continuous traces on the closed cube.
\item \textbf{(Top derivative modulus.)} For every multi-index \(\alpha\) with \(|\alpha|=m\) and every \(x,y\in[-1,1]^d\), the coordinate derivatives, interpreted through their continuous boundary traces, satisfy
\[
|D^\alpha u(x)-D^\alpha u(y)|
\le L_0\left(\sum_{i=1}^d(x_i-y_i)^2\right)^{(\beta-m)/2}.
\]
\item \textbf{(Response bound.)} \(|u(x)|\le M_0\) for every \(x\in[-1,1]^d\).
\end{itemize}
For the affine transport
\[
T(x)=2x-\mathbf1,\qquad x\in[0,1]^d,
\]
the transported function satisfies
\[
u\circ T\in\mathcal H^\beta\!\left(K_{d,\beta}(M_0+L_0)\right)
\]
on \([0,1]^d\), with the intrinsic Hölder ball defined in \cref{obj:def:holder}.
\end{lemmav}

\begin{propositionv}{prop:dorn-a3-free-corollary}[Transfer of Dorn source families]*
Fix parameters satisfying:
\begin{itemize}
\item \textbf{(Dimension and smoothness.)} \(d\in\mathbb N\), \(d\ge1\), and \(\beta>1\).
\item \textbf{(Source constants.)} \(q>3\) and \(M,\sigma_{\min},C,L_0>0\).
\item \textbf{(Adaptation range.)} \(1<\gamma_{\min}<\gamma_{\max}\), with
\[
\Gamma=[\gamma_{\min},\gamma_{\max}].
\]
\end{itemize}

For each \(\gamma\in\Gamma\), let \(\mathcal D_\gamma\) be the maximal envelope of source-law tuples consisting of an observed law \(P\), treated and control mean functions, and a selected propensity \(e_P\), satisfying the following retained restrictions from \citet[Section 2, Assumptions 1--2, pp.~3--7]{Dorn2026GlobalOverlap}: \(X\) is uniform on \([-1,1]^d\); the selected propensity belongs to \((0,1)\) almost surely; conditional on \(X,A\), the outcome is Gaussian with variance \(\sigma_{\min}^2\); both conditional mean functions belong to \(\Sigma^{\mathrm D}(\beta,L_0)\); and
\[
\mathbb E_P\!\left[|Y|^q\mid X,A=1\right]\le M^q,
\qquad
\operatorname{Var}_P(\mu_P(X))\le M,
\qquad
\operatorname{Var}_P(Y\mid X,A)\ge\sigma_{\min}^2,
\]
\[
P\{e_P(X)\le t\}\le Ct^{\gamma-1}
\qquad (0\le t\le1).
\]
Here \(\mu_P\) is the selected treated mean. The source class \(\Sigma^{\mathrm D}(\beta,L_0)\) consists of functions in \(C^m([-1,1]^d)\), with continuous boundary traces of their derivatives, whose derivatives of order
\[
m=\max\{k\in\mathbb Z:k<\beta\}
\]
have \((\beta-m)\)-Hölder seminorm at most \(L_0\), using the source's convention.

Let \(T\), the affine transport from the unit cube to the source cube, and the common response bound \(B_*\) be
\[
T(x)=2x-\mathbf1,
\qquad
B_*=\max\{M,\sigma_{\min}\}.
\]
Define the common transported Hölder radius and the interpolation constant by
\[
L_*=
\sup_{\substack{u\in\Sigma^{\mathrm D}(\beta,L_0)\\
                 \sup_{x\in[-1,1]^d}|u(x)|\le M}}
\|u\circ T\|_{\mathcal H^\beta},
\]
\[
C_{d,\beta}=
\sup_{\substack{M'>0,\ L_0'>0\\
                 u\in\Sigma^{\mathrm D}(\beta,L_0')\\
                 \sup_{x\in[-1,1]^d}|u(x)|\le M'}}
\frac{\|u\circ T\|_{\mathcal H^\beta}}{M'+L_0'}.
\]
The functional \(\|\cdot\|_{\mathcal H^\beta}\) is the full Hölder-radius functional on \([0,1]^d\), including the function, lower derivatives, and top-derivative modulus. Then
\[
0<C_{d,\beta}<\infty,
\qquad
0<L_*\le C_{d,\beta}(M+L_0)<\infty.
\]

Every member of \(\mathcal D_\gamma\), after replacing \(X\) by \(T^{-1}(X)\), belongs to the global-tail causal class \(\mathcal P_\gamma\) with constants
\[
B=B_*,
\qquad L=L_*,
\qquad C=C,
\qquad c_f=1.
\]
In particular, its transported observed law admits a consistency-and-exchangeability causal completion with these constants.

Let \(\widehat\mu_n\) be the equal-cell estimator of \cref{obj:def:estimator}, with smoothness \(\beta\) and response bound \(B_*\), applied to \(n\) independent transported observations. Write \(r_{\gamma,n}\), the oracle risk scale, as
\[
r_{\gamma,n}=n^{-\beta/(2\beta+d\gamma/(\gamma-1))}
\qquad(n\ge1).
\]
There exist \(K>0\) and \(n_0\in\mathbb N\) such that, for every \(n\ge n_0\), every \(\gamma\in\Gamma\), and every source tuple in \(\mathcal D_\gamma\),
\[
\mathbb E_{P^{\otimes n}}
\left[
\sup_{x\in[0,1]^d}
\left|\widehat\mu_n(x)-\mu_P(T(x))\right|
\right]
\le K r_{\gamma,n}.
\]
Every source family \(\mathfrak P\) satisfying the cited Assumptions 1--2 with the fixed displayed constants is nonempty and contained in \(\mathcal D_\gamma\). The same \(K,n_0\) therefore satisfy
\[
\sup_{P\in\mathfrak P}
\mathbb E_{P^{\otimes n}}
\left[
\sup_{x\in[0,1]^d}
\left|\widehat\mu_n(x)-\mu_P(T(x))\right|
\right]
\le K r_{\gamma,n}
\qquad(n\ge n_0).
\]

For every \(\varepsilon>0\), there exist \(c(\varepsilon)>0\) and \(n_0\in\mathbb N\) such that, simultaneously for all \(n\ge n_0\), \(\gamma\in\Gamma\), and source tuples in \(\mathcal D_\gamma\),
\[
P^{\otimes n}\!\left\{
\left\|\widehat\mu_n\circ T^{-1}-\mu_P\right\|_{L^\infty(P)}
>c(\varepsilon)r_{\gamma,n}
\right\}\le\varepsilon.
\]
Here the norm is the essential supremum under the source covariate law, and the estimator is computed from the transported sample.

If \(d\ge2\), then for each \(\gamma\in\Gamma\) there exist positive finite constants \(C',M',\sigma_{\min}',L_0'\) and a source tuple with uniform covariates, Gaussian conditional outcomes, and a selected pointwise propensity representative satisfying all the retained restrictions with these constants, for which
\[
\neg\,
\mathsf A_{3}^{\mathrm{Dorn}}
\bigl(\{P\},(e_P)\bigr),
\]
where the family-level condition is defined in \cref{obj:def:dorn-a3}. The constants in this existence assertion are selected separately for each \(\gamma\).

For any nonempty family \(\mathfrak P\) of source-law tuples sharing a common treated regression curve \(\mu^\circ\), meaning
\[
\mu_P(x)=\mu^\circ(x)
\qquad(P\in\mathfrak P,\ x\in\mathbb R^d),
\]
the minimax pointwise miss-probability value at any evaluation point \(x_0\), sample size \(n\in\mathbb N\), and positive threshold \(cr\), with \(c,r>0\), equals zero.

Finally, for every \(\gamma>1\) and every \(x_0\in[0,1]^d\), there exists \(c_\gamma>0\) such that
\[
c_\gamma r_{\gamma,n}
\le R_{\mathrm{pt}}(n,\gamma,x_0)
\qquad(n\ge1),
\]
where \(R_{\mathrm{pt}}\) is the pointwise minimax expected absolute-error risk of \cref{obj:def:risks} over the causal class with constants \(B_*,L_*,\max\{C,1\}\), and \(c_f=1\). The displayed lower bound is first established for the two-valued potential-outcome subclass risk; monotonicity under class inclusion transfers it to \(R_{\mathrm{pt}}\). The constant \(c_\gamma\) may depend on the evaluation point and the fixed parameters.
\end{propositionv}
% CAUSALSMITH-CITED-SCOPE-BEGIN prop:dorn-a3-free-corollary
\verificationfootnotetext{\textbf{Formalization scope.} This result uses the cited conclusion \citep{Dorn2026GlobalOverlap} (Section 2 Setting and Notation, Assumptions 1--2; author PDF pages 3--7). The cited source proof is not formalized here: Lean verifies its use and all remaining steps. If that cited conclusion is false, the portions of this result that depend on it are not certified.}
% CAUSALSMITH-CITED-SCOPE-END prop:dorn-a3-free-corollary

The completion constant depends on dimension and smoothness alone. Consequently, a common function bound and a common top-derivative modulus yield a common full Hölder radius after affine transport. Applied to the treated regression curves in \cref{obj:prop:dorn-a3-free-corollary}, this statement supplies the finite radius needed by the equal-cell estimator. The continuous boundary traces make the completion applicable on the entire closed cube, including its faces and corners.

\subsection{Source loss and the fixed-family criterion}

For the source comparisons, the regression target is \(\mu_P(x)=\mathbb E_P[Y\mid X=x,A=1]\), with the selected version understood. Let \(\mathcal D_n\) denote the sample of \(n\) independent observed triples, whose sampling law is \(P^{\otimes n}\). For a regression estimator \(\widehat\mu_n\), the source loss is
\[
\|\widehat\mu_n-\mu_P\|_{L^\infty(P)}
=
\operatorname*{ess\,sup}_{x\sim P_X}
|\widehat\mu_n(\mathcal D_n,x)-\mu_P(x)|.
\]
Thus the essential supremum is taken under the source covariate marginal. For a fixed finite overlap exponent \(\gamma_0>1\), write
\[
r^*_{\mu,n}
=
n^{-\beta/(2\beta+d\gamma_0/(\gamma_0-1))}
\qquad(n\ge1)
\]
for the source comparison scale. The fixed-family probability criterion bounds the probability that the source loss exceeds \(c(\varepsilon)r^*_{\mu,n}\), where \(c(\varepsilon)\) is a positive threshold constant allowed to depend on the fixed setup.

A pointwise minimax miss probability also depends on which regression curves the family contains. The next statement records the common-curve obstruction once; later comparisons invoke it by reference.

\begin{lemmav}{lem:arbitrary-family-lower-obstruction}[Known regression minimax obstruction]
Fix parameters satisfying
\begin{itemize}
\item \textbf{(Dimension.)} \(d\in\mathbb N\) and \(d\ge1\).
\item \textbf{(Smoothness.)} \(\beta>0\) and \(L_0>0\).
\item \textbf{(Overlap.)} \(\gamma\in\mathbb R\) and \(\gamma>1\).
\end{itemize}
Let \(\mathfrak P\) be any nonempty source family equipped with selected treated and control regression functions and propensity representatives. Suppose its selected treated regression \(\mu_P\) is the same known function \(\mu^\circ:\mathbb R^d\to\mathbb R\) throughout the family:
\[
\mu_P(x)=\mu^\circ(x)
\qquad
\text{for every }P\in\mathfrak P
\text{ and every }x\in\mathbb R^d.
\]
For a sample size \(n\in\mathbb N\), let \(\mathcal D_n\) denote the observed sample and \(P^{\otimes n}\) its product sampling law. Then, for every evaluation point \(x_0\in\mathbb R^d\), every \(n\in\mathbb N\), and every \(c,r_n>0\), where \(r_n\) is an arbitrary positive comparison scale,
\[
\inf_{T\ \mathrm{measurable}}
\sup_{P\in\mathfrak P}
P^{\otimes n}
\left\{\left|T(\mathcal D_n)-\mu_P(x_0)\right|>c r_n\right\}
=0.
\]
The infimum ranges over all measurable real-valued sample estimators \(T\), and the minimax value is taken in the extended nonnegative reals.

In particular, consider the singleton family whose observed law has mutually independent coordinates
\[
X\sim\operatorname{Unif}([-1,1]^d),
\qquad
A\sim\operatorname{Bernoulli}(1/2),
\qquad
Y\sim N(0,1),
\]
with selected treated and control regression functions identically zero and selected propensity \(e_P(x)=1/2\). There exist common constants \(q>3\), \(M>0\), and \(C>0\) for which this law satisfies all of the following:
\begin{itemize}
\item \textbf{(Selected propensity.)} The selected propensity is a measurable version of the conditional treatment probability and lies in \((0,1)\) almost surely.
\item \textbf{(Moment and variance bounds.)} The treated conditional absolute \(q\)-moment is integrable and satisfies
\[
\mathbb E_P\!\left[|Y|^q\mid X=x,A=1\right]\le M^q
\quad\text{for }P_X\text{-almost every }x,
\qquad
\mathbb E_P\!\left[
\left(\mu_P(X)-\mathbb E_P[\mu_P(X)]\right)^2
\right]\le M.
\]
\item \textbf{(Source smoothness.)} Both selected regression functions belong to the source top-derivative seminorm class \(\Sigma^{\mathrm D}(\beta,L_0)\) specified in \cref{obj:prop:dorn-a3-free-corollary}.
\item \textbf{(Global overlap.)} The selected propensity satisfies
\[
P_X\{e_P(X)\le t\}\le C t^{\gamma-1}
\qquad\text{for every }t\in[0,1].
\]
\item \textbf{(Conditional Gaussian outcomes.)} For \(P_X\)-almost every \(x\), the conditional outcome distribution in each treatment arm is Gaussian with its selected regression mean and variance one.
\item \textbf{(Common local anti-concentration.)} The singleton satisfies
\[
\mathsf A_{3}^{\mathrm{Dorn}}
\bigl(\mathfrak P,(e_P)_{P\in\mathfrak P}\bigr),
\]
with the local ball-ratio inequalities and common positive constants specified in \cref{obj:def:dorn-a3}.
\end{itemize}
At the cube center \(x_0=0\), for every \(n\in\mathbb N\) and every \(c,r_n>0\), this singleton's displayed minimax miss probability equals zero.
\end{lemmav}

The decision underlying \cref{obj:lem:arbitrary-family-lower-obstruction} is the constant sample estimator \(T(\mathcal D_n)=\mu^\circ(x_0)\). Its pointwise error is identically zero throughout the family:
\[
|T(\mathcal D_n)-\mu_P(x_0)|=0
\qquad(P\in\mathfrak P).
\]
The Gaussian singleton shows that the obstruction is compatible with the source moment, smoothness, overlap, and local anti-concentration restrictions.

\subsection{The source upper comparison}

The cited upper comparison concerns an estimator selected for an already fixed family and its common constants. Its probability bound is uniform over members of that family, with the comparison scale determined by the family's fixed overlap exponent. The following proposition retains that quantifier order; its common-regression clause is the application of \cref{obj:lem:arbitrary-family-lower-obstruction} needed to delimit the converse scope.

\begin{propositionv}{lem:dorn-rate-scope}[Fixed-family regression rate scope]*
Fix an integer \(d\ge1\), constants \(\beta,L_0,M,\sigma,C>0\), and \(q>3\). For every fixed finite \(\gamma_0>1\), consider a nonempty source family \(\mathfrak P\) of observed-data laws, equipped with selected treated and control regression curves and selected propensity representatives \((e_P)_{P\in\mathfrak P}\). Suppose that:
\begin{itemize}
\item \textbf{(Design and propensity.)} Under every \(P\in\mathfrak P\), \(X\) is uniform on \([-1,1]^d\), and the selected measurable representative satisfies
\[
e_P(X)=P(A=1\mid X)\in(0,1)
\qquad P\text{-almost surely}.
\]
\item \textbf{(Moment and variance.)} Writing \(\mu_P\) for the selected treated regression curve, every member satisfies
\[
\mathbb E_P\!\left[|Y|^q\mid X,A=1\right]\le M^q,\qquad
\operatorname{Var}_P\!\left(\mu_P(X)\right)\le M,\qquad
\operatorname{Var}_P(Y\mid X,A)\ge\sigma^2,
\]
with the conditional inequalities holding almost surely in each applicable arm.
\item \textbf{(Gaussian regression and smoothness.)} Almost surely with respect to the covariate marginal, the conditional outcome distributions in the treated and control arms are Gaussian with their respective selected regression means and common variance \(\sigma^2\). Both selected regression curves belong to the source seminorm class \(\Sigma^{\mathrm D}(\beta,L_0)\): their derivatives of order
\[
m=\max\{k\in\mathbb Z:k<\beta\}
\]
satisfy Dorn's Hölder-seminorm bound \(L_0\), with exponent \(\beta-m\).
\item \textbf{(Global overlap tail.)} For every \(P\in\mathfrak P\) and \(t\in[0,1]\),
\[
P\{e_P(X)\le t\}\le Ct^{\gamma_0-1}.
\]
\item \textbf{(Family-level local condition.)} On the original cube \([-1,1]^d\), the selected propensity representatives satisfy
\[
\mathsf A_{3}^{\mathrm{Dorn}}
\bigl(\mathfrak P,(e_P)_{P\in\mathfrak P}\bigr)
\]
as defined in \cref{obj:def:dorn-a3}, including its common positive constants \(\rho,\nu,h_0\) and uniform local-ball mass-ratio inequality.
\end{itemize}

Define the fixed-setup comparison scale by
\[
r^*_{\mu,n}
:=n^{-\beta/(2\beta+d\gamma_0/(\gamma_0-1))}.
\]
There exists a regression estimator \(\widehat\mu_n(\mathcal D_n,x)\), where \(\mathcal D_n\) is the sample of \(n\) independent observations, selected after the family, its common constants, and its propensity representatives have been fixed, such that its dependence on the sample is measurable for every \(n\) and \(x\). For every sample size \(n\), member \(P\in\mathfrak P\), and real threshold \(R\), the event
\[
\left\{
\|\widehat\mu_n-\mu_P\|_{L^\infty(P)}
>\max\{Rr^*_{\mu,n},0\}
\right\}
\]
is measurable, where the norm is the essential supremum with respect to the covariate marginal of \(P\). For every \(\varepsilon>0\), there exists a finite constant \(c(\varepsilon)>0\) such that
\[
\limsup_{n\to\infty}\sup_{P\in\mathfrak P}
P^{\otimes n}\!\left\{
\|\widehat\mu_n-\mu_P\|_{L^\infty(P)}
>c(\varepsilon)r^*_{\mu,n}
\right\}
\le\varepsilon,
\]
where \(P^{\otimes n}\) is the product sampling law. The estimator and \(c(\varepsilon)\) may depend on the fixed setup. This upper comparison is the assertion of \citet[Section 2, Assumptions 1--3, Theorem 1(ii)]{Dorn2026GlobalOverlap}.

For every \(\gamma_0>1\), the following two pointwise conclusions also hold. First, let \(\mathfrak P\) be any nonempty source family whose selected treated regression curves equal a common known function \(\mu^\circ\) at every point:
\[
\mu_P(x)=\mu^\circ(x)
\qquad(P\in\mathfrak P,\ x\in\mathbb R^d).
\]
For every evaluation point \(x_0\in\mathbb R^d\), sample size \(n\), and constants \(c,r>0\), its extended nonnegative minimax miss risk is
\[
\inf_{T\ \mathrm{sample\text{-}measurable}}
\sup_{P\in\mathfrak P}
P^{\otimes n}\!\left\{
|T(\mathcal D_n)-\mu_P(x_0)|>cr
\right\}
=0.
\]
Here the infimum ranges over all measurable real-valued functions of the sample.

Second, there exist constants \(q'>3\) and \(M',C'>0\) for which the singleton Gaussian experiment
\[
X\sim\operatorname{Unif}([-1,1]^d),\qquad
A\sim\operatorname{Bernoulli}(1/2),\qquad
Y\sim N(0,1),
\]
with \(X,A,Y\) mutually independent and selected treated mean, control mean, and propensity respectively \(0,0,1/2\), satisfies the displayed design, propensity, moment, variance, Gaussian, smoothness, and global-tail conditions with parameter tuple
\[
(\beta,q',M',1,C',L_0,\gamma_0),
\]
and satisfies the family-level condition of \cref{obj:def:dorn-a3} on the original cube. For this singleton, at \(x_0=0\), every sample size \(n\) and every \(c,r>0\) satisfy
\[
\inf_{T\ \mathrm{sample\text{-}measurable}}
P^{\otimes n}\!\left\{|T(\mathcal D_n)|>cr\right\}
=0.
\]
\end{propositionv}
% CAUSALSMITH-CITED-SCOPE-BEGIN lem:dorn-rate-scope
\verificationfootnotetext{\textbf{Formalization scope.} This result uses the cited conclusion \citep{Dorn2026GlobalOverlap} (Section 2 Notation and Setting, Assumptions 1--3, and Theorem 1 (author PDF pp. 3--7 and 12); Lemma 11 (p. 23); Lemma 13 (p. 24); proof of Theorem 1 (pp. 27--28); Section 5 Conclusion (p. 28)). The cited source proof is not formalized here: Lean verifies its use and all remaining steps. If that cited conclusion is false, the portions of this result that depend on it are not certified.}
% CAUSALSMITH-CITED-SCOPE-END lem:dorn-rate-scope

The source upper assertion in \cref{obj:lem:dorn-rate-scope} is the fixed-setup, uniform-in-\(P\) probability guarantee of \citet[Theorem 1(ii), author PDF p.~12; proof of Theorem 1, pp.~27--28]{Dorn2026GlobalOverlap}. Its displayed scale has no logarithmic factor. The estimator, threshold constant, and overlap exponent remain indexed by the fixed setup, while sample measurability and exceedance-event measurability give the displayed probabilities their ordinary sampling interpretation.

The common-regression clauses give the model-family qualification relevant to the pointwise lower attribution printed in \citet[Theorem 1(i), author PDF p.~12]{Dorn2026GlobalOverlap}. The positive pointwise converse in \cref{obj:prop:dorn-a3-free-corollary} instead concerns the full causal class of \cref{obj:def:risks}, with expected absolute-error loss; its lower bound is supported by the two-valued potential-outcome subclass in \cref{obj:def:bounded-lower-pair,obj:lem:bounded-lower-pair}.

\subsection{The affine transfer}

The transfer component of \cref{obj:prop:dorn-a3-free-corollary} uses the retained source restrictions from \citet[Section 2, Assumptions 1--2, author PDF pp.~3--7]{Dorn2026GlobalOverlap}. The maximal envelopes \(\mathcal D_\gamma\) and the constants \(B_*,L_*\) are defined in that statement. Their role is to place every qualifying source tuple within one response-recovery class after transporting covariates by \(T^{-1}\), where \(T(x)=2x-\mathbf1\).

The ingredients have distinct mathematical functions. Uniform source covariates become uniform covariates on the unit cube, supplying density constant \(c_f=1\). The retained propensity tail controls global overlap at the same exponent. The treated-moment and Gaussian restrictions supply the common response bound \(B_*=\max\{M,\sigma_{\min}\}\), and \cref{obj:lem:fixed-cube-holder-completion} supplies the finite full Hölder radius \(L_*\). The causal-completion clause of \cref{obj:prop:dorn-a3-free-corollary} then identifies the transported observed laws as members of the model in \cref{obj:def:model}.

With this membership established, the equal-cell estimator of \cref{obj:def:estimator} has the expected spatial-supremum guarantee stated in \cref{obj:prop:dorn-a3-free-corollary}, uniformly over the fixed-constant envelopes and the compact overlap range. At \(\gamma=\gamma_0\), its oracle scale \(r_{\gamma,n}\) equals the source scale \(r^*_{\mu,n}\). The accompanying essential-supremum probability bound supplies a comparison in the source loss criterion, with constants common to those envelopes and that exponent range. Every qualifying source subfamily inherits the envelope guarantee through the containment asserted in \cref{obj:prop:dorn-a3-free-corollary}.

\subsection{The Gaussian strictness construction}

The Gaussian existence component of \cref{obj:prop:dorn-a3-free-corollary} concerns the geometric scope of the retained restrictions. For each overlap exponent in the stated range and dimension \(d\ge2\), it supplies a source tuple with uniform covariates, Gaussian conditional outcomes, smooth means, and the required moment and global-tail bounds, together with failure of the selected-version local predicate in \cref{obj:def:dorn-a3}. Its positive finite constants are selected separately for each exponent.

This quantifier structure distinguishes the existence construction from the upper guarantee over a prescribed fixed-constant envelope. The construction establishes that the retained Gaussian restrictions accommodate the selected-propensity geometry described by the corollary. The affine-transfer component supplies recovery guarantees for each envelope with its displayed common constants, while \cref{obj:lem:dorn-rate-scope} supplies the preprint comparator for fixed families satisfying the additional local condition. Together, these components identify the smoothness completion, family scope, and treatment geometry underlying \cref{obj:prop:dorn-a3-free-corollary}.
